\section{Hướng dẫn tái lập demo}
\label{sec:repro}

Demo được thiết kế để \emph{bất kỳ ai cũng chạy lại được}, đúng tiêu chí
\emph{reproducibility} của BetterBench~\cite{betterbench}. Có ba cơ chế
bảo đảm tính tất định.

\subsection{Ba trụ cột tái lập}

\begin{enumerate}
  \item \textbf{Dữ liệu cố định theo seed}: tập con benchmark lấy bằng
    \texttt{SEED = 0} và ghi ra \texttt{data/*.jsonl} $\Rightarrow$ mọi người nhận
    \emph{đúng cùng item, cùng thứ tự}.
  \item \textbf{Sinh văn bản greedy}: \texttt{do\_sample=False} $\Rightarrow$ đầu
    ra không phụ thuộc ngẫu nhiên; chạy lại cho đúng cùng văn bản.
  \item \textbf{Cache theo nội dung}: mỗi lời gọi lưu theo
    \texttt{sha256(model\_id, prompt)} trong \texttt{results/predictions/}. Cache
    lưu văn bản thô, điểm \emph{tính lại} mỗi lần $\Rightarrow$ tái sinh toàn bộ
    bảng/biểu đồ \emph{không cần GPU}, và thêm model mới chỉ tốn chi phí cho phần
    mới.
\end{enumerate}

\subsection{Cách A --- Kaggle}

\begin{enumerate}
  \item Mở \texttt{notebook/2\_run\_all.ipynb}. Settings $\to$ Accelerator
    $\to$ \texttt{GPU T4}; Internet $\to$ On.
  \item Add-ons $\to$ Secrets $\to$ thêm \texttt{HF\_TOKEN} (HF read token). Chấp
    nhận licence Gemma trên Hugging Face (các model khác ungated).
  \item \emph{Run all}. Vì greedy + cache, bấm lại ra \emph{đúng cùng số}. Lần đầu
    $\sim$1.5--2 giờ; các lần sau gần như tức thì.
\end{enumerate}

\subsection{Cách B --- Máy có GPU NVIDIA $\geq 8$GB}

\begin{verbatim}
python -m venv .venv && source .venv/bin/activate
pip install -r requirements.txt
huggingface-cli login          # cần đã accept licence Gemma
python src/evaluate.py         # -> results/per_item.csv
python src/analysis.py         # -> results/*.csv + images/*.png
\end{verbatim}

\subsection{Cách C --- Không có GPU (chỉ kiểm tra pipeline)}

\begin{verbatim}
pip install pandas matplotlib numpy
MINIBENCH_SYNTHETIC=1 python src/evaluate.py --fake
MINIBENCH_SYNTHETIC=1 python src/analysis.py
\end{verbatim}

Đường này dùng \texttt{FakeRunner} + dữ liệu tổng hợp để chứng minh \emph{toàn bộ
luồng} chạy được trên máy không GPU; \textbf{số liệu sinh ra không dùng để báo cáo}
--- chỉ để kiểm tra code.

\subsection{Sản phẩm tái lập}

Sau khi chạy, \texttt{results/} chứa: \texttt{per\_item.csv} (từng câu),
\texttt{accuracy.csv}, \texttt{ranking.csv}, \texttt{coverage.csv},
\texttt{mmlu\_by\_subject.csv}, \texttt{metric\_gap.csv} (+\texttt{\_examples}),
\texttt{robustness.csv}, \texttt{error\_analysis.csv}, và \texttt{images/*.png}.
Mã nguồn và cache prediction được công khai trong repo, nên kết quả trong báo cáo
này có thể được \emph{tái kiểm} đầy đủ.

\subsection{Giới hạn của tái lập}

Tái lập \emph{tất định} chỉ đúng với cùng phiên bản thư viện (\texttt{transformers},
\texttt{bitsandbytes}) và cùng phần cứng; lượng tử hoá 4-bit và kernel GPU khác nhau
có thể gây sai khác nhỏ ở biên. Tuy nhiên cơ chế cache văn bản thô giúp các \emph{con
số trong báo cáo} luôn tái sinh được từ chính các prediction đã lưu, độc lập với
phần cứng.
